\section{Título y Objetivos}
\label{sec:titulo/objetivos}

\noindent \textbf{Práctica 1:} Determinación de la densidad de un cilindro macizo por distintos métodos e instrumentos, y su expresión con sentido físico.

\noindent \textbf{Práctica 2:} Determinación experimental de la aceleración de la gravedad ($g$) mediante el estudio del péndulo simple ideal, aplicando regresión lineal sobre la relación $T^2$ en función de $L$.

\noindent Los objetivos específicos de ambas prácticas son:
\begin{itemize}
    \item Realizar mediciones directas con distintos instrumentos (calibre, regla milimetrada, probeta graduada, balanza electrónica, regla para longitud del hilo).
    \item Calcular magnitudes indirectas (volumen, densidad, período, aceleración gravitacional) propagando correctamente las incertezas.
    \item Comparar resultados obtenidos por diferentes métodos y analizar la compatibilidad de sus intervalos de incerteza.
    \item Aplicar regresión lineal para obtener el valor de $g$ y sus incertezas (error de la pendiente y de la ordenada al origen).
    \item Expresar todos los resultados con sentido físico (valor representativo $\pm$ incerteza absoluta, con unidades).
\end{itemize}
